\section{Introduction}

Chat models built on large language models (LLMs) are now ordinary software components: they answer questions, transform text, and produce structured output that other programs parse. The Hugging Face Hub distributes hundreds of open chat models, many of them small enough to run on a laptop, and exposes them both as downloadable weights and through a hosted inference API. A developer who embeds such a model faces the same question as with any third-party component: on which inputs does it misbehave? Answering this question is hard for three reasons. Chat models rarely crash, so the crash oracle that drives most fuzzers reveals little. Their input language is natural language, which has no parser whose structure a fuzzer could exploit. And for most prompts there is no single correct output, which makes the test-oracle problem acute \citep{barr2015oracle}. Hand-written examples, in turn, explore only a small and biased part of the input space.

Fuzzing research offers tools for each of these obstacles. Grammar-based fuzzing generates large numbers of structured, valid inputs from a grammar \citep{holler2012langfuzz}, and coverage of the grammar itself can guide generation when the internals of the system under test cannot be observed \citep{havrikov2019kpath}. Differential testing replaces a missing specification by comparing several implementations on the same input \citep{mckeeman1998differential}, and metamorphic testing replaces it by comparing a system's outputs on related inputs \citep{chen1998metamorphic}. Neither comparison, however, establishes that an output is wrong: two LLMs may legitimately phrase the same answer differently, and a disagreement does not say which model erred. Existing LLM fuzzers mostly target jailbreaks, whose oracle is whether an attack succeeded \citep{yu2023gptfuzzer}, and metamorphic frameworks for LLMs compare each model only with itself \citep{hyun2024metal,cho2025metamorphic}. We did not find work that generates prompts from a grammar, guides generation by grammar coverage, and applies both a cross-model differential oracle and a cross-paraphrase metamorphic oracle to open chat models, nor work that measures how often these oracles are right when the true answer is known.

\textbf{Research problem and objective.} This paper studies whether grammar-based fuzzing combined with differential and metamorphic testing is an effective way to reveal functional failures in small open chat models, and what limits the conclusions it supports. We make the oracle problem tractable by design: our grammar generates \emph{verifiable} prompts, whose expected answer follows from the values drawn during generation (arithmetic, letter counting, string reversal, JSON output, and verbatim echo of hard Unicode text), plus a category of free-form requests with no expected answer. Because the ground truth is known, we can measure what the differential and metamorphic oracles would have reported without it. Specifically, we aim to (i)~quantify and characterize the failures the approach reveals, (ii)~measure the recall and precision of the differential and metamorphic oracles against ground truth, (iii)~evaluate whether coverage-guided expansion improves on random expansion, and (iv)~estimate how often the automatic oracles misjudge a response.

This paper makes the following contributions:
\begin{enumerate}
  \item \code{chatfuzz}, an open-source grammar-based, coverage-guided fuzzer for chat models that generates paraphrase families with built-in ground truth and supports local Hugging Face models and the Hugging Face Inference Providers API.
  \item An empirical study on three small open chat models (Qwen2.5-0.5B-Instruct, SmolLM2-360M-Instruct, and TinyLlama-1.1B-Chat) that characterizes 271 oracle failures in 420 graded responses, including failures that only paraphrasing reveals.
  \item A quantitative comparison of oracles: using ground truth as reference, we report how much of the failure set the differential and metamorphic oracles recover, how often they raise false alarms, and how rarely all models agree on the same wrong answer.
  \item An evaluation of coverage guidance on grammar coverage and failure diversity, and a manual validation of the automatic oracles that locates their main source of error.
\end{enumerate}

Section~\ref{sec:related} reviews related work, Section~\ref{sec:rqs} states the research questions, Sections~\ref{sec:method} and~\ref{sec:setup} describe the method and the experimental setup, Section~\ref{sec:results} reports the results, Section~\ref{sec:discussion} discusses them, Section~\ref{sec:threats} covers threats to validity, and Section~\ref{sec:conclusion} concludes.
